%! Author = Saeid Yazdani
%! Date = 9/9/2024


\section{Conclusion and Future Work}
\label{sec:conclusion-and-future-work}

\subsection{Conclusion}

The scope of this work was to demonstrate the applicability of a hybrid
digital-analog approach for power supply design with a dynamic compensation
scheme for the testing of semiconductors in \acs{ATE} systems.

Advancements in FPGA logic size and speed have enabled the implementation of
real-time adaptive control schemes.
These advancements provide the means to address the challenges presented by
modern semiconductor devices, particularly in ensuring high-quality power
delivery during testing.

The primary focus of this work was the optimization of signal waveforms at
the pins of \acs{DUT}s, where anomalies such as ringing, overshoot, and
undershoot in the supply rails can significantly degrade test quality and yield.
These anomalies could damage the \acs{DUT} in worst-case scenarios.
This issue is critical because, although modern devices may survive short
overvoltage pulses such as ringing, their long-term reliability can be
compromised due to the potential latent failures induced during the test phase
and only show up in the field.

As a result, the quality of power rail waveforms is of utmost importance.
Fast rise and fall times without ringing are essential to avoid any
potential degradation in device behavior.
This work demonstrated that safe and predictable waveforms can be achieved
through adaptive feedback control, where parameter sets are switched dynamically
— similar to the adaptive positioning systems used in fast robotics.

One of the key challenges discussed in this research was the limitation
of current-generation \acs{ATE} \acs{DPS}, which are often used in
ganged-mode to provide the necessary currents for demanding \acs{DUT}s.
This mode comes with its own set of challenges.
A single power supply that can provide the necessary current at all common
voltage levels would be the more flexible solution as proposed in this work.

A prototype system was developed to validate the advantages of this hybrid
approach.
Challenges and shortcomings encountered in the proposed
three-stage driver circuit were studied and solutions were provided.

Measurements demonstrated that modern \acs{FPGA}, \acs{ADC}, and \acs{DAC}
components allow for rapid responses, which are critical for maintaining
the stability and precision of the power delivery system.

The primary limitation observed was the amplitude-discrete nature of digital
signals.
To overcome this, cascaded converters with coarse and vernier stages were
developed, achieving the necessary amplitude resolution for precise control.

In conclusion, this work has successfully demonstrated the feasibility and
advantages of a hybrid digital-analog approach in \acs{ATE} systems, providing
a pathway for future improvements in semiconductor testing and ensuring more
robust, reliable device behavior during and after the test phase.

\subsection{Future Work}


\noindent\textbf{Upgrading the Digital Block:}

Between the time of starting this research work and the time of writing this
text, the \acs{FPGA} technology has seen technological leaps in several key
areas driven by the need for higher performance and improved design flexibility.
The \acs{HVS} can benefit from new generation of \acs{FPGA} chips
in the following areas:

\begin{itemize}[topsep=8pt,itemsep=4pt,partopsep=4pt, parsep=4pt]
    \item \textbf{Integration of Artificial Intelligence:}
    FPGAs have significantly evolved to support \ac{AI} and \ac{ML} workloads.
    Many FPGA families now include dedicated AI engines, such as \ac{TPU} or
    other neural network accelerators, built into the fabric.
    These advancements make FPGAs highly competitive for edge AI and machine
    learning applications that can benefit the \acs{HVS}'s adaptive
    compensation scheme.

    \item \textbf{SoC Integration:}
    FPGAs have become highly integrated heterogeneous platforms and feature
    hardware based embedded processors (e.g., ARM Cortex cores).
    The \acs{HVS} can significantly benefit from the speed up of
    communication between the operator software and the FPGA,
    lower design complexity and smaller footprint on the \ac{PCB}
    by replacing the discrete \acs{MCU} by an integrated one.

    \item \textbf{High-Speed \acs{I/O}:}
    Modern \acs{FPGA}s support much higher data rates, often exceeding
    \SI{100}{\giga\bit\per\second} for transceivers, and they offer built-in
    support for the latest connectivity standards such as 6th generation
    \acs{PCIe} and fast \acs{DDR}5.
    A high \ac{I/O} speed of the \acs{FPGA} in \acs{HVS} can aalow for faster
    data transactions.
    Including high-speed external memory can extend the acquisition memory
    used for both control loop and datapoint storage.

\end{itemize}

\pagebreak
\noindent\textbf{Upgrading the Analog Block:}

Just like \acs{FPGA}s, there are new and faster power amplifiers,
\acs{ADC}s and \acs{DAC}s available on the market.
For the current booster stage, upgrading from silicon-based \acs{MOSFET}s to
wide bandgap \acs{MOSFET}s (\emph{Silicon Carbide} or \emph{Gallium Nitride})
can be beneficial in achieving faster response rates.
\emph{Silicon Carbide} offers lower on-resistance
and reduced switching losses, which increases efficiency at higher frequencies.
Faster switching speeds and high thermal conductivity of such \acs{MOSFET}s
are essential factors to consider.
These new technologies should be evaluated in future iterations of the
\acs{HVS}.

\noindent\textbf{Synchronization of Multiple \acs{HVS} Units:}

Exact synchronization between \acs{DPS} units is mandatory
for modern devices with multiple voltage rail requirements (e.g., \SI{5}{\volt}
for analog, \SI{3.3}{\volt} for I/O, and \SI{1.3}{\volt}
for core supplies) as otherwise there is a high risk of latch-up.
To mitigate this, a fixed startup sequence for all individual supply rails is
necessary.
As the proposed system voltages are monitored digitally, this setup phase can
be easily programmed between multiple units of the proposed system to ensure
proper sequencing and avoid related failures.

\noindent\textbf{Integration in Recognized ATE Systems:}

The \acs{HVS} prototype was developed and tested as a stand-alone
\acs{DPS}/\acs{SMU}.
Integration of this system in industry-standard \ac{ATE} systems requires
direct collaboration with \acs{ATE} manufacturers.
The current iteration of the \acs{HVS} offers a simple but effective communication
protocol, and library files for easy integration are well prepared.
Another option to expose the \acs{HVS} to other hardware is the integration of
the \acs{HVS}'s \acs{API} as an extension to the already available
\emph{SpecScribe} \autocite{markert, schott}, a specification tool for analog
hardware developed at TU Chemnitz.

\noindent\textbf{Moving to a Single Board Design:}

The current iteration of the \acs{HVS} hardware uses the backplane-based
prototyping approach. The driver stage, analog frontend and control boards
are the most critical sections of the hardware where any signal distortion can
affect the overall behavior of the system.
The signals travel through multiple connectors and the backplane to
reach their destination.
Because of this, the communication rates between system components are not
at their maximum as the priority in the prototype was
the proof of functionality.
By combining critical sections of the hardware in one board problems can
be avoided and maximum transfer rates can be achieved.



